% Paso 4: Guía de Investigación Histórica
% Proyecto: Análisis Exegético Profundo
% Ministerio "La Gloria es del Señor"

\documentclass[11pt,a4paper]{article}
\usepackage[utf8]{inputenc}
\usepackage[spanish]{babel}
\usepackage[margin=2.5cm]{geometry}
\usepackage{graphicx}
\usepackage{fancyhdr}
\usepackage{xcolor}
\usepackage{tcolorbox}
\usepackage{enumitem}
\usepackage{tabularx}
\usepackage{multirow}
\usepackage{amsfonts}
\usepackage{fontspec}
\usepackage{titlesec}

% Configuración de colores
\definecolor{forestgreen}{RGB}{27,67,50}
\definecolor{sagegreen}{RGB}{45,90,39}
\definecolor{olivegreen}{RGB}{64,145,108}
\definecolor{mintgreen}{RGB}{82,183,136}
\definecolor{lightmint}{RGB}{116,198,157}
\definecolor{softcream}{RGB}{216,243,220}
\definecolor{papergray}{RGB}{118,128,150}

% Configuración de encabezados
\pagestyle{fancy}
\fancyhf{}
\fancyhead[L]{\textcolor{forestgreen}{\textbf{Análisis Exegético Profundo}}}
\fancyhead[R]{\textcolor{papergray}{Paso 4: Contexto Histórico-Cultural}}
\fancyfoot[C]{\textcolor{papergray}{\thepage}}
\fancyfoot[R]{\textcolor{papergray}{\scriptsize Ministerio "La Gloria es del Señor"}}

% Configuración de títulos
\titleformat{\section}
{\color{forestgreen}\Large\bfseries}
{\thesection}{1em}{}

\titleformat{\subsection}
{\color{sagegreen}\large\bfseries}
{\thesubsection}{1em}{}

% Configuración de cajas
\tcbuselibrary{skins}
\newtcolorbox{infobox}[1]{
    colback=softcream,
    colframe=olivegreen,
    fonttitle=\bfseries,
    title=#1,
    left=10pt,
    right=10pt,
    top=10pt,
    bottom=10pt
}

\newtcolorbox{checklistbox}{
    colback=white,
    colframe=mintgreen,
    left=15pt,
    right=15pt,
    top=10pt,
    bottom=10pt,
    boxrule=2pt
}

\newtcolorbox{historybox}{
    colback=lightmint!20,
    colframe=sagegreen,
    left=10pt,
    right=10pt,
    top=8pt,
    bottom=8pt,
    boxrule=1.5pt
}

\newtcolorbox{historytimeline}{
    colback=olivegreen!10,
    colframe=olivegreen,
    left=10pt,
    right=10pt,
    top=8pt,
    bottom=8pt,
    boxrule=1.5pt
}

\begin{document}

% Título principal
\begin{center}
    {\Huge\color{forestgreen}\textbf{PASO 4}}\\[0.5cm]
    {\Large\color{sagegreen}\textbf{CONTEXTO HISTÓRICO-CULTURAL}}\\[0.3cm]
    {\large\color{papergray}Investigación del Trasfondo del Mundo Bíblico}\\[1cm]
    
    \begin{tcolorbox}[colback=olivegreen!10, colframe=olivegreen, width=0.8\textwidth]
        \centering
        \textcolor{forestgreen}{\textbf{Proyecto: Análisis Exegético Profundo}}\\
        \textcolor{papergray}{Metodología de 8 Pasos para Estudio Bíblico Riguroso}
    \end{tcolorbox}
\end{center}

\vspace{1cm}

% Información del estudiante
\begin{infobox}{Información del Proyecto}
\begin{tabularx}{\textwidth}{|X|X|}
\hline
\textbf{Nombre del Estudiante:} & \hspace{6cm} \\[0.8cm]
\hline
\textbf{Texto Bíblico:} & \hspace{6cm} \\[0.8cm]
\hline
\textbf{Libro Bíblico:} & \hspace{3cm} \quad \textbf{Autor:} \hspace{3cm} \\[0.8cm]
\hline
\textbf{Fecha Aproximada:} & \hspace{6cm} \\[0.8cm]
\hline
\end{tabularx}
\end{infobox}

\section{Objetivo del Paso 4: Contexto Histórico-Cultural}

La investigación histórico-cultural nos permite entender el mundo del autor y de la audiencia original. Este trasfondo es esencial para una interpretación precisa, ya que las Escrituras fueron escritas en contextos específicos con realidades culturales particulares.

\subsection{Descripción del Proceso}

Investigaremos el trasfondo histórico, geográfico, cultural y social del texto para comprender el mundo en que vivían el autor y sus lectores originales. Este contexto ilumina el significado pretendido del texto.

\begin{infobox}{Recursos Recomendados}
\textbf{Historia:} Bright, Merrill, Hoehner\\
\textbf{Cultura:} deSilva, Bailey, Pilch \& Malina\\
\textbf{Geografía:} Atlases bíblicos, Rasmussen\\
\textbf{Arqueología:} Hoerth, Currid, Hoffmeier
\end{infobox}

\section{Lista de Verificación}

\begin{checklistbox}
\textbf{Marca cada elemento completado:}

\begin{itemize}[leftmargin=1cm]
    \item[$\square$] Investigar el período histórico específico
    \item[$\square$] Estudiar la geografía y ambiente físico
    \item[$\square$] Analizar las circunstancias del autor
    \item[$\square$] Entender la audiencia original y su situación
    \item[$\square$] Considerar trasfondos religiosos relevantes
    \item[$\square$] Examinar el contexto político y social
    \item[$\square$] Investigar costumbres y prácticas culturales
    \item[$\square$] Consultar evidencia arqueológica pertinente
\end{itemize}
\end{checklistbox}

\newpage

\section{Contexto Histórico}

\subsection{Cronología y Período Histórico}

\begin{historytimeline}
\textbf{Fecha del Evento/Escrito:} \hspace{8cm}\\[0.5cm]
\textbf{Período Histórico:} \hspace{8cm}\\[0.5cm]
\textbf{Dinastía/Imperio Dominante:} \hspace{8cm}
\end{historytimeline}

\textbf{Eventos históricos significativos del período:}
\begin{enumerate}[leftmargin=1cm]
    \item \hspace{12cm}
    \vspace{0.3cm}
    \item \hspace{12cm}
    \vspace{0.3cm}
    \item \hspace{12cm}
    \vspace{0.3cm}
    \item \hspace{12cm}
    \vspace{0.3cm}
    \item \hspace{12cm}
\end{enumerate}

\subsection{Línea de Tiempo Relevante}

\begin{historybox}
\textbf{Construye una línea de tiempo de 50-100 años centrada en tu texto:}
\end{historybox}

\vspace{4cm}
\begin{center}
[Espacio para línea de tiempo - incluye fechas clave, eventos y personajes]
\end{center}
\vspace{2cm}

\subsection{Situación Política}

\textbf{¿Quién gobernaba en el momento del texto?}
\vspace{2cm}

\textbf{¿Cuál era la situación política de Israel/la iglesia?}
\vspace{2cm}

\textbf{¿Qué tensiones o conflictos políticos existían?}
\vspace{2cm}

\newpage

\section{Contexto Geográfico}

\subsection{Ubicación Geográfica}

\begin{historybox}
\textbf{Lugar(es) Mencionado(s):} \hspace{8cm}\\[0.5cm]
\textbf{Región/Provincia:} \hspace{8cm}\\[0.5cm]
\textbf{Características Topográficas:} \hspace{8cm}
\end{historybox}

\textbf{Descripción del ambiente físico:}
\vspace{2cm}

\textbf{¿Cómo afectaba la geografía la vida diaria?}
\vspace{2cm}

\textbf{¿Hay significado simbólico en la geografía mencionada?}
\vspace{2cm}

\subsection{Rutas de Comercio y Comunicación}

\textbf{¿Estaba el lugar en rutas comerciales importantes?}
\vspace{1.5cm}

\textbf{¿Cómo afectaba esto el intercambio cultural?}
\vspace{1.5cm}

\section{Contexto Cultural y Social}

\subsection{Estructura Social}

\begin{tabularx}{\textwidth}{|X|X|}
\hline
\textbf{Clases Sociales Presentes} & \textbf{Roles y Características} \\
\hline
& \\[1.5cm]
\hline
& \\[1.5cm]
\hline
& \\[1.5cm]
\hline
& \\[1.5cm]
\hline
\end{tabularx}

\subsection{Costumbres y Prácticas Culturales}

\textbf{Costumbres familiares relevantes:}
\vspace{2cm}

\textbf{Prácticas comerciales o económicas:}
\vspace{2cm}

\textbf{Tradiciones religiosas pertinentes:}
\vspace{2cm}

\textbf{Roles de género en la sociedad:}
\vspace{2cm}

\newpage

\section{Contexto Religioso}

\subsection{Trasfondo del Antiguo Testamento (si aplica)}

\textbf{¿Qué antecedentes del AT son relevantes?}
\vspace{2cm}

\textbf{¿Hay alusiones o citas del AT en el texto?}
\vspace{2cm}

\textbf{¿Qué expectativas mesiánicas/escatológicas existían?}
\vspace{2cm}

\subsection{Situación Religiosa Contemporánea}

\begin{tabularx}{\textwidth}{|X|X|X|}
\hline
\textbf{Grupo Religioso} & \textbf{Creencias Principales} & \textbf{Relevancia al Texto} \\
\hline
Fariseos & & \\[1cm]
\hline
Saduceos & & \\[1cm]
\hline
Esenios & & \\[1cm]
\hline
Zelotas & & \\[1cm]
\hline
Otros: & & \\[1cm]
\hline
\end{tabularx}

\section{Autor y Audiencia}

\subsection{Circunstancias del Autor}

\textbf{¿Cuál era la situación personal del autor?}
\vspace{2cm}

\textbf{¿Qué motivó la escritura de este texto?}
\vspace{2cm}

\textbf{¿Dónde estaba el autor cuando escribió?}
\vspace{2cm}

\subsection{Audiencia Original}

\begin{historybox}
\textbf{Audiencia Primaria:} \hspace{8cm}\\[0.5cm]
\textbf{Ubicación Geográfica:} \hspace{8cm}\\[0.5cm]
\textbf{Composición Étnica:} \hspace{8cm}\\[0.5cm]
\textbf{Situación Socioeconómica:} \hspace{8cm}
\end{historybox}

\textbf{¿Qué desafíos enfrentaba la audiencia?}
\vspace{2cm}

\textbf{¿Qué conocimiento previo tenían del tema?}
\vspace{2cm}

\textbf{¿Cómo habría sido recibido originalmente el mensaje?}
\vspace{2cm}

\newpage

\section{Evidencia Arqueológica}

\subsection{Descubrimientos Relevantes}

\begin{tabularx}{\textwidth}{|X|X|X|}
\hline
\textbf{Descubrimiento} & \textbf{Fecha/Lugar} & \textbf{Relevancia al Texto} \\
\hline
& & \\[1.5cm]
\hline
& & \\[1.5cm]
\hline
& & \\[1.5cm]
\hline
\end{tabularx}

\subsection{Cultura Material}

\textbf{¿Qué objetos, edificios o prácticas mencionadas en el texto han sido confirmados arqueológicamente?}
\vspace{3cm}

\section{Síntesis del Contexto}

\subsection{Factores Contextuales Más Importantes}

\textbf{¿Cuáles son los 3 factores histórico-culturales más importantes para entender tu texto?}

\begin{enumerate}[leftmargin=1cm]
    \item \hspace{12cm}
    \vspace{1cm}
    \item \hspace{12cm}
    \vspace{1cm}
    \item \hspace{12cm}
    \vspace{1cm}
\end{enumerate}

\subsection{Implicaciones para la Interpretación}

\textbf{¿Cómo afecta este trasfondo histórico-cultural tu comprensión del texto?}
\vspace{3cm}

\textbf{¿Qué aspectos del mensaje original podrían perderse sin este contexto?}
\vspace{3cm}

\textbf{¿Hay algún malentendido común que este contexto ayuda a corregir?}
\vspace{3cm}

\begin{infobox}{Principio Hermenéutico}
El contexto histórico-cultural no determina el significado, pero sí lo delimita. Busca entender el mundo del texto sin proyectar nuestra cultura moderna sobre él.
\end{infobox}

\section{Siguiente Paso}

Con la investigación histórico-cultural completada, procederás al \textbf{Paso 5: Género Literario}, donde analizaremos las características literarias del texto.

\vspace{1cm}

\begin{center}
\begin{tcolorbox}[colback=mintgreen!20, colframe=mintgreen, width=0.7\textwidth]
\centering
\textcolor{forestgreen}{\textbf{¡Investigación Completa!}}\\
\textcolor{papergray}{Has explorado el rico trasfondo histórico-cultural del texto.}
\end{tcolorbox}
\end{center}

\end{document}